\section{Reproducing and using the artifacts}
\label{app:reproduce}
The package contains the article sources and compiled PDF, the three producers,
a search-independent checker, exhaustive finite oracles, unit tests, recorded
inputs and receipts, example JSON files, a small uncompiled Lean lemma specimen,
and build/replay instructions. No third-party source tree or font file is
included. The Python tools require Python 3.10 or newer; only 3.13.5 was used
for the recorded run.

From the package root:
\begin{lstlisting}
python -S tests/test_core.py
python -S prototype/replay.py
python -S prototype/run_experiments.py --output reproduced-results

python -S prototype/forge_worker.py search vass \
  examples/mutex.model.json --output candidate.json
python -S prototype/forge_worker.py check vass \
  examples/mutex.model.json --certificate candidate.json

python -S prototype/forge_worker.py search nominal \
  examples/lru_fifo.model.json
python -S prototype/forge_worker.py search mixed \
  examples/owner_pool.model.json
\end{lstlisting}
The search CLI validates inputs, runs the producer, and checks any non-unknown
result before printing it. Exit code zero means an accepted model result,
including a valid negative witness; exit code two means unknown; exit code one
means an input/checking error. An accepted result is still a Python-model
result, not a Lean proof. The check-only path does not import producers.

The default experiment output directory is \code{reproduced-results}; it does
not overwrite the retained evidence unless explicitly directed to do so.
The recorded seed determines the generated corpus, but timings depend on the
runtime and host. Article compilation uses pdfLaTeX with standard packages and
no shell escape:
\begin{lstlisting}
cd article
pdflatex -interaction=nonstopmode -halt-on-error forge-infinite-state.tex
pdflatex -interaction=nonstopmode -halt-on-error forge-infinite-state.tex
pdflatex -interaction=nonstopmode -halt-on-error forge-infinite-state.tex
\end{lstlisting}

\section{Model and certificate schema notes}
\label{app:schema}
\paragraph{Counter model.}
The fields are \code{dimension}, \code{controls}, \code{initials},
\code{transitions}, and \code{bad}. An initial family contains \code{control},
\code{base}, and a list of \code{rays}. A transition contains \code{src},
\code{dst}, \code{consume}, and \code{produce}; a bad threshold contains
\code{control} and \code{vector}. All entries have validated dimensions;
negative integers and Python Booleans are not accepted as natural numbers.

\paragraph{Counter receipts.}
A positive object has \code{schema: 1}, \code{kind: "safe"}, and one
\code{basis} list per control. A negative object has \code{kind: "unsafe"},
an \code{initial} family index, \code{parameters}, and original
\code{transitions} indices. Debug \code{stats} are not evidence. The axes
receipt's mathematical payload is:
\begin{lstlisting}
{"schema": 1, "kind": "safe", "basis": [[[1, 1]]]}
\end{lstlisting}

\paragraph{Equality models.}
A nominal input file is a pair of machine objects. Each supplies register,
constant, and control counts; tags; initial control and anchored register values;
and ordered rules. Selectors are \code{["input"]}, \code{["null"]},
\code{["reg", i]}, and \code{["const", j]}. Guards are Boolean values or
\code{eq}, \code{not}, \code{and}, and \code{or} syntax. Updates list one
selector per register. Outputs carry a \code{bool} or \code{atom} tag and the
appropriate expression. Fixed constants use labels $0,\ldots,c-1$.

\paragraph{Equality receipts.}
A positive object has \code{kind: "equivalent"} and a \code{states} list of
\code{[leftControl, rightControl, jointRegisters]}. The checker requires distinct
orbits, not necessarily canonical labels. A negative object has
\code{kind: "different"} and a concrete \code{word} of \code{[tag, atom]}
entries. It need not stop at the first mismatch to be valid; shortestness is a
producer property tested separately, not part of receipt soundness.

\paragraph{Mixed models and receipts.}
A mixed model combines a single register-machine header with counter initial
families and nondeterministic rules carrying \code{consume}/\code{produce}.
Bad clauses carry input-independent register guards. A safe receipt has
\code{states} entries \code{[control, registers]} and one counter basis per
entry. An unsafe receipt has initial-family parameters and concrete
\code{steps} triples \code{[ruleId, tag, atom]}. The receipt does not include
a trusted compiled edge table.

\paragraph{Unknown.}
Search may return \code{kind: "unknown"} with a reason and statistics.
The certificate checkers intentionally do not accept it as a proof. The JSON
schema is a prototype interchange format; extra diagnostic fields are ignored,
while the mathematical fields and model dimensions are checked. It is not a
versioned public Forge protocol yet.

\section{Source audit and implementation status}
\label{app:status}
Full baseline identifiers are:
\begin{lstlisting}
Forge  58ea206bd7ad501b930add568151217bcc7f82a2
Leant  6bf05ad78c467989e68290f2d08bbed40802d485
Djex   e8778f4ebd63e1f9b9b410fa4de8d14a8a04c9e5
\end{lstlisting}
The inspected Forge README gives a Lean v4.34.0 baseline and Mathlib
\code{1cf325a0cf67}; the full Mathlib identifier is
\begin{lstlisting}
1cf325a0cf67aca2b04d76b5380ff6a9e410aefa
\end{lstlisting} Leant's reported native
Djex dependency is \code{e237e866}; that dependency identifier must not be
silently replaced by the separate Djex repository-head identifier above.

The main Forge audit paths were:
\begin{lstlisting}
README.md
article/sections/06-closure.tex
article/sections/15-conclusion.tex
article/sections/B-provenance.tex
lean/Forge/Design/Contracts.lean
lean/Forge/Design/Runtime.lean
lean/Forge/Closure/Covers.lean
\end{lstlisting}
The Leant current commit response included
the complete indexing-composition diagnostic report; Djex's README supplied its
checked synthesis scope and embedding entry points. Their archived materials
are not reproduced in this package.

\begin{longtable}{@{}>{\raggedright\arraybackslash}p{.38\textwidth}>{\raggedright\arraybackslash}p{.20\textwidth}>{\raggedright\arraybackslash}p{.36\textwidth}@{}}
\toprule
Artifact or assertion & Status & Evidence \\
\midrule\endhead
Counter, nominal, mixed producers & Executed & Stored corpus and generator. \\
Independent Python receipt replay & Executed & 650 accepted objects; no producer
imports in replay. \\
Differential finite-oracle checks & Executed & 640 agreements, subject to described
generator and evaluator sharing. \\
Exact fragment soundness and completeness & Mathematical arguments & Proofs in
Sections~\ref{sec:antichain}--\ref{sec:product}; not formalized here. \\
Scalar/vector Lean specimen & Uncompiled & Source included, no \code{sorry};
\code{NOT\_RUN}, no claimed axiom output. \\
Lean certificate checker and source frontend & Proposed & Module boundaries,
proof obligations, and acceptance gates only. \\
Leant/Djex semantic-summary integration & Proposed & Candidate-flow design; original
indexing and universe queries remain untouched. \\
Performance versus Lean tactics & Not measured & No end-to-end tactic comparison
or claimed solved-goal gain. \\
\bottomrule
\end{longtable}

\section{Concluding assessment}
The meaningful extension is a change in what Forge can finitely represent:
upward regions of unbounded natural states, equality orbits of states containing
arbitrary names, and the exact product of these representations under an
explicitly restricted syntax. The proofs identify both the reason for finite
certificates and the cases where those certificates cease to be complete.

The prototypes do more than demonstrate plausible examples. They discover
positive receipts, produce executable failures, preserve explanations under
subsumption and renaming, compare against exhaustive finite oracles on suitable
fragments, and replay stored evidence without importing searches. They also
leave the important boundary visible: a model theorem becomes a Lean source
theorem only through a verified checker and a proved translation. Building that
narrow path is the recommended next contribution to Forge.
